\documentclass[conference]{IEEEtran}
% NOTE: IEEEtran.cls must be available in your TeX distribution.
% On this system it was NOT found at setup time. Install via:
%   sudo apt-get install texlive-publishers
% or place IEEEtran.cls alongside this file.

\usepackage{cite}
\usepackage{amsmath,amssymb,amsfonts}
\usepackage{algorithmic}
\usepackage{graphicx}
\usepackage{textcomp}
\usepackage{xcolor}
\usepackage{hyperref}

\begin{document}

\title{6G-Enabled Hybrid Quantum-Secure Architecture for Electronic Health Record Sharing}

\author{
\IEEEauthorblockN{Ramana Sree K V}
\IEEEauthorblockA{\textit{TODO: Department, Institution} \\
\textit{TODO: City, Country} \\
TODO: email@example.com}
\and
\IEEEauthorblockN{Verona Ann Mariya}
\IEEEauthorblockA{\textit{TODO: Department, Institution} \\
\textit{TODO: City, Country} \\
TODO: email@example.com}
}

\maketitle

\begin{abstract}
% TODO: Write abstract last, after Results and Conclusion are finalized.
% TODO: 150-250 words summarizing problem, gap, proposed architecture,
% methodology, and (once available) key results.
\end{abstract}

\begin{IEEEkeywords}
% TODO: e.g., 6G, Post-Quantum Cryptography, Quantum Key Distribution,
% Electronic Health Records, Hybrid Security Architecture
\end{IEEEkeywords}

\section{Introduction}
% TODO: Motivation, problem statement, contributions, paper organization.
% Do NOT insert claims or results until verified against actual sources
% and experiments.

\section{Background and Related Work}
% TODO: Survey of 6G security, PQC, QKD, EHR security, edge-assisted
% approaches. Populate from research/literature_matrix.csv once literature
% collection begins.

\section{Research Gap}
% TODO: Derive directly from research/research_gaps.md. Do not state a gap
% here that is not documented and justified there.

\section{Proposed Architecture}
% TODO: Describe the hybrid quantum-secure 6G architecture for EHR sharing.
% This section must be clearly marked as PROPOSED / conceptual until and
% unless it is implemented and validated in simulation.

\section{Methodology}
% TODO: Describe research design, simulation tools, evaluation metrics,
% and assumptions. To be filled in only once tools/methods are decided.

\section{Experimental Setup}
% TODO: Simulation environment, parameters, datasets (if any), baseline
% comparisons. No experiments have been run yet.

\section{Results and Discussion}
% TODO: Populate ONLY with actual results generated from experiments/
% and simulation/. Do not fabricate figures, tables, or numbers.

\section{Security Analysis}
% TODO: Threat model, security proofs/arguments, comparison against
% classical and quantum adversary models.

\section{Limitations}
% TODO: Honest discussion of scope, assumptions, and constraints of the
% proposed architecture and evaluation.

\section{Conclusion}
% TODO: Summarize contributions and future work. Write only after all
% preceding sections are finalized.

% \section*{Acknowledgment}
% TODO (optional)

\bibliographystyle{IEEEtran}
\bibliography{../references/references}

\end{document}
